% Generated draft: requires mathematical and source review.
\documentclass[11pt]{article}
\usepackage[margin=25mm]{geometry}
\usepackage{fontspec}
\setmainfont{Latin Modern Roman}
\usepackage{amsmath,amssymb,mathtools,unicode-math}
\setmathfont{Latin Modern Math}
\usepackage{enumitem,needspace,xurl,hyperref}
\setlength{\parindent}{0pt}
\setlength{\parskip}{5pt}
\setlength{\emergencystretch}{3em}
\newcommand{\N}{\mathbb{N}}
\newcommand{\Z}{\mathbb{Z}}
\newcommand{\Q}{\mathbb{Q}}
\newcommand{\R}{\mathbb{R}}
\newcommand{\C}{\mathbb{C}}
\newcommand{\F}{\mathbb{F}}
\newcommand{\A}{\mathbb{A}}
\newcommand{\PP}{\mathbb{P}}
\newcommand{\E}{\mathbb{E}}
\newcommand{\eps}{\varepsilon}
\newcommand{\dd}{\,\mathrm{d}}
\begin{document}
\section*{UnsolvedMath ID: 701153}
\textbf{Source problem code:} catalogue entry
\subsection{Definitions and mathematical statement}
Let \(F\) be a finite field and let \(A\in\operatorname{Mat}_d(F)\). Write \(m_A(t)\) and \(\chi_A(t)\) for its minimal and characteristic polynomials. Following the source, call \(A\) \emph{separable} if \(m_A(t)\) has no repeated roots in an algebraic closure of \(F\), and call \(A\) \emph{cyclic} if \(m_A(t)=\chi_A(t)\). For a finite matrix group \(G\), define
\[
s(G)=\frac{|\{A\in G:m_A\text{ has no repeated roots}\}|}{|G|},\qquad
c(G)=\frac{|\{A\in G:m_A=\chi_A\}|}{|G|}.
\]
For a classical matrix group series denoted in the source by \(\mathrm{X}(d,q)\), with dimension parameter \(d\) and finite-field parameter \(q\), consider
\[
S(\mathrm{X};q)=\lim_{d\to\infty}s(\mathrm{X}(d,q)),\qquad
C(\mathrm{X};q)=\lim_{d\to\infty}c(\mathrm{X}(d,q)).
\]
Here \(q\) is fixed when taking a dimension limit, and dimensions must run through those admissible for the chosen series. The source introduces these limits without separately stating their existence hypotheses.

\textbf{Question (P. M. Neumann, 2002).} For the unitary, symplectic, and orthogonal classical group series, are \(S(\mathrm{X};q)\) and \(C(\mathrm{X};q)\) rational functions of \(q\)? In the usual interpretation of rationality, this asks, for each specified series \(\mathrm{X}\), whether there are \(R_{\mathrm{X},s}(T),R_{\mathrm{X},c}(T)\in\Q(T)\) such that
\[
S(\mathrm{X};q)=R_{\mathrm{X},s}(q),\qquad
C(\mathrm{X};q)=R_{\mathrm{X},c}(q)
\]
for every admissible prime power \(q\), with the functions defined at those values.

\textbf{Unresolved qualifications.} The supplied record does not specify the exact unitary and orthogonal matrix groups, the unitary field-parameter convention, or how orthogonal parity and form type are handled in the dimension limits. Consequently these choices, and the existence of the corresponding limits, require clarification or confirmation before the question determines a unique collection of precise assertions. The interpretation using \(\Q(T)\) above makes the usual meaning of rationality explicit; the source does not specify the coefficient field.

The source attributes rationality of the corresponding general linear group limits to independent work of G. E. Wall and J. Fulman. Its registry status is \emph{open}. A separate background entry reports an explicitly unverified AI claim of a negative answer; this draft records that qualification without treating the claim as an established result.
\subsection{Short English statement}
For each unitary, symplectic, or orthogonal matrix group family, fix the finite-field parameter and let the dimension grow. Do the proportions of matrices with squarefree minimal polynomial, and of matrices whose minimal and characteristic polynomials coincide, converge to rational functions of the field parameter? The exact group conventions and dimension sequences need clarification. The supplied negative-answer claim is explicitly unverified.
\subsection{Sources}
\begin{itemize}
\item Local UnsolvedMath registry, permanent ID \(701153\), exact source problem code \texttt{catalogue entry}, title \emph{Kourovka Notebook Problem 15.65}; supplied complete source record, upstream revision \texttt{d371ee917ac9647fbb06cd6fc529deec6f3dda82}.
\item P. M. Neumann, Problem 15.65, \emph{The Kourovka Notebook (Unsolved Problems in Group Theory)}, 15th issue (2002); also identified in the supplied record with the 21st issue edition (2026), arXiv:1401.0300v48. The collection description names E. I. Khukhro and V. D. Mazurov as editors.
\item G. E. Wall, \emph{Bull. Austral. Math. Soc.} \textbf{60}, no. 2 (1999), 253--284. The supplied statement attributes evaluation and rationality of the general linear group limits to this work; no article title is supplied.
\item J. Fulman, \emph{J. Group Theory} \textbf{2}, no. 3 (1999), 251--289. The supplied statement attributes independent evaluation and rationality of the general linear group limits to this work; no article title is supplied.
\item The supplied background and status review report an \emph{unverified AI claim} of a negative answer, attributed there to M. Kinyon and J. Urban, \emph{Preprint} (2026): \url{https://kourovkanotebookorg.wordpress.com/wp-content/uploads/2026/09/kourovka-mathematics.pdf}. The review describes this as marked with \(\#\) in the October 2026 21st issue, arXiv:1401.0300v48, and classifies the evidence as \texttt{unrefereed\_claim}.
\item The supplied status review is dated \texttt{2026-10-16} and retains canonical status \texttt{open}. This is supplied dated curation, not an independently verified literature assessment. No cited source has been independently consulted for this draft.
\end{itemize}
\end{document}
